% =============================================================================
\chapter{Pipeline Selector}
\label{ch:pipeline-selector}
% =============================================================================

The pipeline selector is the landing page for the app shell. Everyone
starts here after login: pick a pipeline to enter its workspace, or
continue in global mode by navigating via the sidebar.

% --------------------------------------------------------------------------
\section{PipelineSelectorPage}
\label{sec:pipeline-selector-page}
% Route: /
% Component: app/pages/PipelineSelectorPage.tsx

\screenshot{03-pipeline-selector-hero.png}{Pipeline selector — card grid of available pipelines.}

\begin{itemize}
  \pagelement{Header}{Page title and a short hint about what a pipeline is.}
  \pagelement{Card grid}{One card per pipeline (code, name, status, owner).}
  \pagelement{Pipeline card}{Click navigates into the pipeline workspace (Chapter~\ref{ch:pipeline-workspace-x121}).}
  \pagelement{Create pipeline action}{Visible for admins; opens the admin flow.}
  \pagelement{Empty state}{Rendered when no pipelines exist yet, with a call-to-action.}
  \pagelement{Error state}{Rendered when the backend call fails.}
\end{itemize}

\screenshot{03-pipeline-selector-empty.png}{Empty state — no pipelines yet.}
\screenshotnarrow{03-pipeline-selector-card-detail.png}{A single pipeline card showing code, name and status.}

\begin{designnote}
The selector is intentionally the app root (\route{/}) so a fresh
session always surfaces pipeline choice, even for users who have been
working inside a pipeline previously.
\end{designnote}
